\documentclass[11pt,a4paper]{article}
\usepackage[margin=25mm]{geometry}
\usepackage{fontspec}
\usepackage{microtype}
\usepackage{amsmath,amssymb}
\usepackage{booktabs,longtable,array,tabularx}
\usepackage{enumitem}
\usepackage[numbers,sort&compress]{natbib}
\usepackage{xcolor}
\usepackage{hyperref}
\hypersetup{colorlinks=true,linkcolor=blue!45!black,citecolor=blue!45!black,urlcolor=blue!45!black,
  pdftitle={Humanoid Social Navigation: Literature Review and Research Opportunities for Geoff-Sim}}
\urlstyle{same}
\setlength{\emergencystretch}{3em}
\setlength{\parskip}{0.35em}
\setlist{nosep,topsep=0.4em}
\newcolumntype{P}[1]{>{\raggedright\arraybackslash}p{#1}}
\newcolumntype{Y}{>{\raggedright\arraybackslash}X}
\newcommand{\repo}[1]{\path{#1}}
\newcommand{\gap}[1]{\paragraph{Research implication.} #1}
\title{Humanoid Social Navigation:\\State of the Art, Evidence Limitations,\\and Research Opportunities for Geoff-Sim}
\author{Project literature review and research-positioning document}
\date{Literature search and repository inspection: 22 September 2026}

\begin{document}
\maketitle
\begin{abstract}
Geoff-Sim integrates Unitree G1 locomotion, MuJoCo, ROS~2, online mapping,
pedestrian tracking, human-aware planning, collision supervision, and experimental
benchmarking. This review evaluates 50 identity-verified research publications
and preprints spanning 1995--2026. It connects classical collision avoidance,
learning-based crowd navigation, trajectory prediction, legged control,
onboard perception, and social-navigation evaluation. Recent work already
demonstrates LiDAR-conditioned humanoid avoidance, cross-embodiment navigation,
and physically simulated multi-task benchmarks. Consequently, the strongest
positioning is not the first integration of walking and social navigation.
It is a falsifiable investigation of how calibrated locomotion response,
observation age, uncertainty, and planning/supervision interaction affect
closed-loop collision risk and mission liveness on a walking humanoid.
The repository has unusually useful information-source and experiment-validity
instrumentation, but its documented onboard trials still include collisions
and timeouts. Five research opportunities are ranked, with explicit baselines,
ablations, metrics, and claim boundaries. Source identity verification is
separated from abstract-level screening, selected full-text analysis, and
experimental replication; the latter was not undertaken.
\end{abstract}
\tableofcontents
\newpage

\section{Executive assessment}
\label{sec:executive}
\textbf{Recommended core research question:}
\emph{Can a calibrated, observation-age-aware model of a humanoid's executable
motion improve the collision--completion trade-off of social navigation under
partial observability, relative to matched geometric, temporal, and fixed-margin
baselines?}

This question fits the current implementation closely. It makes the robot's
actual delayed response, rather than its requested planar velocity, the object
of prediction. It also recognizes that a planner using a fresh message can
still be acting on an old observation. Its novelty must be established by
the specific model, calibration procedure, decision rule, and controlled
evidence, because dynamics-aware avoidance and biped crowd navigation are
established topics \citep{dwa,vo,biped2016}.

Three conclusions organize the review:
\begin{enumerate}
\item \textbf{The field has moved beyond isolated obstacle avoidance.}
Recent systems combine learned locomotion with perception and navigation;
the closest humanoid paper maps LiDAR and proprioception directly to motor
commands, and FLUX reports transfer to a physical G1
\citep{safehumanoid,flux}. A general claim that previous research ignores
humanoids, dynamics, or raw sensing is untenable.
\item \textbf{The remaining problem is measurable closed-loop reliability.}
Forecast accuracy, nominal command feasibility, or a visually successful
demonstration does not establish low collision risk over complete missions.
Perception loss, stale observations, tracking association, command shaping,
and stopping transients can interact. Current repository documentation
provides concrete examples of this gap.
\item \textbf{Evaluation can be a substantive contribution if it tests a new
scientific question.} Separating oracle control, sensor-only control, and
independent ground-truth scoring enables causal diagnosis. The contribution
would be a validated protocol and new findings about failure mechanisms,
not simply another simulator or another collection of ROS nodes.
\end{enumerate}

The most defensible near-term paper is a simulation-based methods and evaluation
study. Hardware transfer and human comfort should be treated as subsequent
validation targets, not inferred from MuJoCo clearance measurements.

\section{Scope, search procedure, and verification standard}
\label{sec:method}
\subsection{Review type and research questions}
This is a \textbf{structured critical scoping review}, not an exhaustive
systematic review or a meta-analysis. It asks:
\begin{enumerate}
\item What is already established about planning around people while respecting
the dynamics and sensing constraints of a mobile or legged robot?
\item Which recent humanoid and cross-embodiment systems are direct competitors?
\item What evidence distinguishes geometric avoidance, social acceptability,
and reliable mission-level autonomy?
\item Which gaps can this repository realistically investigate, and what would
falsify the proposed advantages?
\end{enumerate}

The search cutoff is 22 September 2026. Foundational work is retained when it
defines a necessary baseline or assumption. The latest included preprint,
DPed-VLN, was submitted on 18 September 2026 \citep{dped}; its inclusion is an
abstract-level horizon scan, not a claim of mature validation. Date and venue
are taken from retrieved records, not inferred from identifier prefixes.

\subsection{Search and inclusion procedure}
Discovery used Crossref title/bibliographic searches, arXiv title and topic
searches, and references encountered in closely related papers. OpenAlex was
used as an auxiliary locator for an accessible copy of the older biped paper.
Representative topic strings were ``social humanoid navigation'', ``socially
aware robot navigation'', ``safe navigation biped crowd'', ``LiDAR legged
collision avoidance'', ``humanoid sim-to-real'', and exact method titles.
The reproducible query families and screening decisions are recorded in
\repo{SEARCH_LOG.md}. Exact selected identifiers are in \repo{sources.json}.

Inclusion required a direct connection to at least one of: human-aware
navigation; crowd perception/prediction; executable collision avoidance;
legged/humanoid control; simulation/middleware used by the project; or
evaluation methodology. Foundational infrastructure papers are included to
separate enabling technology from methodological novelty. Generic language
models, manipulation-only systems, and unrelated humanoid demonstrations
were not used to inflate the count. Title-search mismatches and duplicate
versions were rejected. For example, a search for FAST-LIO2 returned both an
SSRN record and the IEEE article; the IEEE record was selected.

The retained corpus contains \textbf{50 distinct works: 41 Crossref DOI records
and nine arXiv records}. An arXiv citation does not imply that the work is
unpublished; it means the cited version was verified through arXiv rather than
an independently checked proceedings record. Publication-year differences
between preprints, early access, and journal issues are normal. This review
uses the year returned for the selected record.

\subsection{What ``verified'' means here}
\label{sec:verification}
Verification has three distinct levels:
\begin{description}
\item[Identity verification (all 50).] A primary bibliographic record was
retrieved from Crossref or arXiv. Title, author list, year, DOI/arXiv identifier,
and source URL were extracted and checked. The evidence snapshot records the
retrieval timestamp and SHA-256 of the response. This establishes that a
matching paper record exists; it does not establish the truth of its findings.
\item[Content inspection.] Registry abstracts and additional arXiv abstracts
were inspected where available. The closest recent competitors---ASAP,
Safe and Comfortable Locomotion, FLUX, and OmniNavBench---received selected
full-text methods/evaluation analysis. Specific table-level criticisms below
refer to those inspected versions. Other works primarily supply conceptual
context and candidate baselines; uninspected numerical claims are avoided.
\item[Replication.] No external algorithm was reimplemented or rerun for this
review. No new robot experiment or human-subject evaluation was performed.
Repository results are explicitly attributed to existing development records.
\end{description}

\repo{verified_sources.json} retains the identity evidence and available
abstracts. \repo{READING_AUDIT.md} identifies the inspected material and remaining
access limitations. \repo{verify_sources.py} checks the local snapshot,
bibliography, and citation keys offline, and can refresh the online records.
Full-text access was uneven: the HAL copy of the 2016 biped paper presented an
access challenge, so this review does not attribute detailed experimental
claims to it. Search ranking, indexing delays, language coverage, and a
single-pass synthesis limit completeness. A final publication should extend
forward/backward citation screening and independently review its novelty claim.

\section{The current work and its evidence boundary}
\label{sec:project}
\subsection{Architecture relevant to the review}
Geoff-Sim is a simulation-first Unitree G1 research platform. Its documented
pipeline connects MuJoCo/DDS robot state and simulated sensors to ROS~2,
SLAM Toolbox, Nav2, dynamic-person tracking, short-horizon prediction,
human-aware command selection, a collision supervisor, and a learned
locomotion controller. This is an appropriate architecture for investigating
the interfaces between autonomy layers, rather than only isolated planner
performance.

The repository inspection covered \repo{README.md},
\repo{docs/ONBOARD_SOCIAL_NAVIGATION.md},
\repo{docs/DYNAMIC_ROUTE_PLANNING.md},
\repo{docs/BENCHMARK_EXPERIMENTS.md}, and the implementation
\repo{motion_envelope.py} in the bridge package. File hashes and the repository
revision/dirty-state context are retained in \repo{project_evidence.json}.
The working tree contained ongoing development; the reviewed files, not the
base commit alone, define this assessment's project snapshot.

\subsection{Oracle and sensor-only modes are different experiments}
The oracle social-waypoint mode uses simulator-derived person tracks and a
known static map. The onboard mode uses local-frame raycast LiDAR, IMU and
joint information, scan-registration odometry, online SLAM, observation-based
tracking, and constant-velocity pedestrian forecasts. Evaluation truth remains
available for scoring but is excluded from the documented navigation path.
The source audit is a graph/configuration check rather than a security
sandbox. This distinction is scientifically important: privilege can enter
through localization, goal selection, maps, or tracks, not only through a
planner's direct subscriptions.

The onboard detector is geometric, not a semantic camera or learned person
classifier. The odometry front end assumes approximately planar walking,
and its covariance is not a calibrated localization uncertainty model.
Consequently, the present system provides an experimental platform for
perception-aware navigation, rather than a completed solution to open-world
person recognition or general uneven-terrain localization.

\subsection{What the recorded trials establish}
The onboard status document explicitly states that an earlier oracle 24/24
result has not been reliably reproduced with onboard sensing. Development
trials include an eight-goal completion and several collisions and timeouts.
These runs span implementation changes and must not be pooled into a
single algorithm success rate. The dynamic-route validation records
136 route updates but only three of four goals and one collision.
This establishes active integration, not reliable avoidance.

Timing records are similarly informative but conditional. One onboard trial
reports planner p95 47.115\,ms, tracker p95 53.501\,ms, and track age at output
p95 0.305\,s. A later route-validation trial reports request-to-matching-route
p95 2.649\,s, including retries and blocked-route periods. These quantities
measure different things: node computation time, observation age, and delivered
route latency cannot be substituted for one another. None is a hard real-time
guarantee. The reported 45\,ms candidate-search budget can overrun.

\subsection{The key technical opening}
The shared motion envelope already normalizes command bounds and models
acceleration-limited, delayed planar motion. Its source explicitly states that
the defaults are \emph{conservative engineering assumptions, not a fitted G1
response certificate}. That admission identifies a tractable research step:
measure the command-to-motion relationship and test whether a calibrated
model changes safety decisions. Current velocity caps vary by profile, so
experiments must freeze the policy, gains, normalization, delay model, and
active command limits instead of quoting one cap as universal.

The batch tools already preserve failed/invalid/missing trials and distinguish
control-source evidence from scoring truth. These are useful foundations,
but their presence does not by itself validate the response model or prove
a novel navigation method.

\section{Classical and learning-based social navigation}
\label{sec:navigation}
\subsection{Geometric and dynamically constrained avoidance}
The Dynamic Window Approach (DWA) searches feasible velocities subject to
the robot's dynamics and stopping considerations \citep{dwa}. Velocity
obstacles (VO) reason about relative motion and select velocities outside
collision-producing regions; the original work also intersects these with
acceleration-admissible velocities \citep{vo}. Therefore, merely adding an
acceleration limit to a geometric planner is not a sufficient novelty claim.
Their continuing value is as interpretable baselines: failures can be traced
to motion estimates, admissible sets, or objective functions.

Timed trajectory optimization offers another classical comparator. Sparse
trajectory optimization underlying the timed-elastic-band family jointly
adjusts a path and its temporal structure \citep{teb}. It makes time and
kinematic constraints explicit, but practical behavior still depends on
initialization, obstacle representations, and the fidelity of its executable
motion model. A G1 can have command deadbands, axis coupling, turning-arc
restrictions, and balance-related transients that are not captured by a
generic planar robot model.

ORCA distributes collision-avoidance responsibility between interacting
agents \citep{orca}. This is powerful when agents conform sufficiently to
the reciprocal assumptions. The repository's independently walking actors
do not react to the robot, so an unmodified reciprocal assumption can be
misleading. An ORCA comparison should declare whether people participate
in the same avoidance rule, remain nonreactive, or follow recorded
trajectories. A planner baseline and a pedestrian simulator are different
uses of ORCA and should not be conflated.

\gap{The useful distinction is not ``dynamics-aware versus all previous
kinematic work''. It is whether a particular model of the deployed
locomotion policy, command filtering, and observation delay improves
decisions beyond established feasible-velocity and timed-trajectory methods.}

\subsection{Interaction, reciprocity, and freezing}
Social-force modeling provides a compact account of interactions among
pedestrians \citep{socialforce}. Its interpretability makes it useful for
simulation and baseline reasoning. It does not make a simulated crowd a
validated model of all human behavior, particularly for groups, distraction,
or unexpected encounters with a humanoid. Crowd generation should therefore
be an experimental factor rather than an unquestioned ground truth.

Trautman and colleagues explicitly study cooperation in dense crowds
\citep{freezing}. The central conceptual challenge is that treating every
possible human motion as an obstacle can prevent progress, while assuming
cooperation can be unsafe when the assumption is wrong. The same tension
appears in Geoff-Sim: stronger margins may prevent a local encounter but
produce waypoint timeouts; a rapid command can improve progress while
increasing collision or fall risk. Performance should be evaluated on a
collision--completion frontier rather than a single ``safety score''.

The early human-aware-navigation survey and subsequent taxonomies emphasize
that obstacle avoidance is only one part of navigation among people
\citep{survey2013,survey2024}. A robot must also make its behavior legible,
respect interaction context, and avoid disrupting people. A successful
nonreactive-actor benchmark validates unilateral avoidance under a defined
motion model. It does not measure cooperative negotiation or human comfort.

\subsection{Learned interaction models}
CADRL learns a value function for decentralized collision avoidance, moving
much of the interaction computation into offline training \citep{cadrl}.
SARL uses attention to reason about collective human interactions rather
than only independent robot--human pairs \citep{sarl}. Relational Graph
Learning adds graph-structured representations and predictive lookahead
\citep{rgl}. DS-RNN incorporates spatial and temporal relationships and
addresses partial observability; its abstract also reports transfer to a
TurtleBot~2i \citep{dsrnn}. These methods provide a useful progression from
pairwise local reasoning to temporal crowd representations.

The critical comparison is the information and dynamics contract. A policy
consuming tracked positions and velocities does not perform the detector's
task. A benchmark using a circular planar robot does not necessarily test
balance, upper-body swing, or actual stopping transients. Conversely, it
would be inaccurate to label all learning-based methods ``oracle-only''
or ``simulation-only'', particularly given DS-RNN and newer sensor-level
systems. Each comparator needs its own documented observation interface.

The 2025 review and benchmarking paper identifies strengths of learning-based
and end-to-end approaches in coordination scenarios \citep{survey2023}.
This supports including a learned navigation comparator if resources permit.
However, importing a policy trained with different pedestrian reactions,
speed ranges, sensors, or reward conventions can create a distribution-shift
test rather than a fair algorithm comparison. Such transfer is valuable,
but should be labeled and analyzed separately from matched training.

\gap{A modular response-aware approach is attractive for interpretability
and data efficiency, not because learned policies cannot avoid people.
Its case depends on comparable input information and a demonstrated
advantage under delay, missed detections, or changed locomotion response.}

\section{Prediction and onboard perception}
\label{sec:perception}
\subsection{From point forecasts to interaction-conditioned distributions}
Social LSTM introduced learned interaction-aware trajectory prediction;
Social GAN explicitly represents diverse socially plausible futures;
Trajectron++ combines graph structure, heterogeneous agent data, and
dynamical feasibility \citep{sociallstm,socialgan,trajectron}. Trajectron++
can condition predictions on an ego-agent plan, which is especially relevant
when robot action influences pedestrian behavior. Rudenko et al.'s survey
provides a broad taxonomy of trajectory modeling and contextual information
\citep{prediction}.

These developments challenge a single constant-velocity forecast over long
horizons, but do not imply that a larger network will improve this project's
closed-loop behavior. Short-horizon constant velocity is a valuable baseline:
it is inexpensive, auditable, and can be more stable than acceleration
extrapolation from noisy surface centroids. The onboard profile's use of
constant velocity is therefore defensible as a starting point. Its limits
must be measured at crossings, turns, stops, and occlusion exits.

Prediction error averaged over entire trajectories can conceal the errors
that cause a collision. Multi-sample best-of-$K$ displacement error can also
reward producing one good sample without providing calibrated probabilities.
For planning, the relevant questions are whether the occupied region covers
the person at encounter time, whether probability estimates are calibrated,
and whether the chosen robot motion remains feasible. This is a synthesis
of the planning problem, not a claim that all reviewed predictors omit
these issues.

For a track observed at $t_o$ and used at $t_c$, define its age
$a=t_c-t_o$. A minimal time-aligned mean prediction is
\begin{equation}
\hat p_h(t_c+\tau)=\hat p_h(t_o)+\hat v_h(t_o)(a+\tau).
\label{eq:age}
\end{equation}
The corresponding occupied region must also widen with measurement,
velocity, and unmodeled-motion uncertainty. Equation~\ref{eq:age} is only
a proposed baseline formulation; it is not evidence that the current
implementation has calibrated coverage. Repeated publication of an old
track must not reset $t_o$.

\subsection{Detection and tracking are not solved by geometry alone}
DROW establishes learned detection from sparse two-dimensional range data,
initially emphasizing wheelchairs \citep{drow}. DR-SPAAM uses spatial
attention and recurrent feature updates for person detection in range scans
\citep{drspaam}. These are pertinent perception baselines, but transferring
their assumptions to a moving torso-mounted 3D LiDAR requires an explicit
projection and domain-adaptation study. A flat scan through legs is not
equivalent to clipped, oscillating torso returns.

JRDB provides egocentric human detection/tracking data in built environments
\citep{jrdb}. SCAND provides human-demonstrated socially compliant navigation
on morphologically different platforms, including a quadruped and a wheeled
robot \citep{scand}. The former is useful for characterizing missed detections
and association; the latter is useful for behavioral imitation and judging
which maneuvers humans demonstrate. Neither automatically supplies the
G1-specific LiDAR mounting, sway pattern, or command response of this project.

The repository already attempts partial-body track retention and measured
fragment recovery. These are practical responses to a genuine problem:
an observed surface centroid is not a stable estimate of a person's body
center. Differentiating that centroid may produce an apparent acceleration
that belongs to visibility geometry rather than human motion. Ground-truth
evaluation should separate center bias, velocity error, identity continuity,
track age, and complete misses. The current nearest-track diagnostic is
explicitly not identity-matched detection recall; it should not be reported
as such.

\subsection{Localization, mapping, and the moving sensor}
FAST-LIO2 demonstrates direct LiDAR--inertial estimation, while ORB-SLAM3
provides a visual/visual--inertial alternative \citep{fastlio,orbslam}.
They identify strong adjacent baselines and underline that state estimation
is a developed field. The project's planar scan-registration front end
should be positioned as an engineering component unless a new estimation
method is separately evaluated.

Walking creates a feedback loop between body motion, sensing, and planning.
Attitude errors can affect ground removal and scan projection; localization
drift can shift both maps and tracks; dynamic-person masking can remove useful
geometry. Correctly handling unknown map cells does not by itself resolve
these effects. Conversely, adding a more capable estimator may not fix
collision failures caused by stale tracks or unmodeled braking.

\gap{A controlled oracle-to-onboard decomposition is stronger than a single
end-to-end success number. Independently substitute estimated pose, observed
tracks, and online maps to locate which uncertainty source changes decisions.
Then test interactions between those sources rather than assuming their
effects add linearly.}

\section{Legged control and the closest humanoid competitors}
\label{sec:humanoid}
\subsection{Locomotion competence versus navigation competence}
Massively parallel training reduces the time needed to learn locomotion
policies \citep{rudin}. RMA learns rapid adaptation to changing locomotion
conditions \citep{rma}. Robust perceptive locomotion combines exteroception
and proprioception to handle difficult terrain \citep{rough}.
These advances make a learned low-level policy an increasingly practical
component, but their task objectives differ from mission-level navigation
through a crowd. Good velocity tracking in a locomotion benchmark does not
specify a safe pedestrian-avoidance trajectory.

DeepMimic establishes example-guided physics-based motion imitation
\citep{deepmimic}; ExBody transfers expressive whole-body behavior to
humanoids while relaxing lower-body imitation to support velocity tracking
\citep{expressive}. These are relevant to natural motion and potentially
legibility, but motion resemblance is not evidence of perceived social
comfort. Human evaluation would be needed to establish that relationship.

The boundary between locomotion and navigation is already an active research
area. ANYmal Parkour uses a hierarchy that accounts for locomotion skills
when choosing how to traverse obstacles \citep{parkour}. The wheeled-legged
navigation work integrates adaptive locomotion, mobility-aware local
navigation, and large-scale planning, reporting kilometer-scale missions
\citep{wheeledleg}. Earlier work explicitly addresses safe biped navigation
in crowds \citep{biped2016}. These are direct counterexamples to an overly
broad claim that navigation has always ignored legged execution.

\subsection{ASAP: why nominal simulator dynamics are not enough}
ASAP pretrains motion-tracking policies, collects real-world trajectories,
learns a residual action model to align simulated response, and fine-tunes
the policy in the corrected simulator \citep{asap}. The inspected version
compares transfer between simulators and to a real G1. Its tables examine
both open-loop dynamics mismatch and closed-loop motion imitation.
This separation is important: fitting transitions and improving a deployed
policy are related but distinct achievements.

For Geoff-Sim, ASAP strengthens the motivation to calibrate response rather
than trust nominal gains. It also limits the novelty of ``learning dynamics
corrections''. The proposed contribution here would operate at a different
interface: a short-horizon reachable motion envelope for a frozen locomotion
policy, used to select or veto navigation commands. It would not duplicate
ASAP's whole-body policy fine-tuning objective. Dynamics randomization is
another established transfer tool \citep{domain}, but randomized training
does not supply a calibrated deployment-time stopping bound by itself.

ASAP is not a crowd-navigation evaluation. Its imitation errors and skill
success rates cannot be numerically compared with pedestrian collision or
waypoint completion rates. The useful transfer is methodological: evaluate
model fidelity and downstream behavior separately, on held-out conditions.

\subsection{LiDAR-conditioned safe and comfortable humanoid locomotion}
The most directly overlapping 2025 work is \emph{End-to-End Humanoid Robot
Safe and Comfortable Locomotion Policy} \citep{safehumanoid}. It combines
LiDAR features and proprioceptive history, uses a constrained Markov decision
process, translates discrete barrier-function ideas into training costs,
and incorporates comfort-oriented rewards. The inspected paper reports
Unitree G1/Livox Mid-360 deployment and onboard execution. Its actor/critic
description explicitly distinguishes privileged training information from
deployment inputs.

This prior work makes ``first sensor-driven safe humanoid avoidance'' an
indefensible positioning. It also offers a useful contrast between an
end-to-end learned controller and a modular planner plus supervisor.
However, several details constrain how strongly its results can be used:
\begin{itemize}
\item Table~III reports 30 trials per scenario and rounded success values of
86\% for the proposed method versus 70\% for P3O in the dynamic-agent case.
The same proposed method reports 60\% in a narrow passage. These are
scenario-dependent outcomes, not universal avoidance guarantees.
\item Table~II reports nonzero unsafe-region time even for the proposed
method. The training discussion uses a 0.8\,m safety distance, while that
table defines unsafe space below 0.6\,m. The evaluation thresholds must be
kept separate rather than merged into a single certified clearance claim.
\item The described barrier condition is converted into an optimization
cost during learning, rather than enforced as an online action constraint.
A theoretical invariant-set condition is not automatically inherited by
an approximately trained policy under sensing and dynamics errors.
\item Comfort is operationalized through reward terms and geometric/dynamic
proxies. The inspected experiments do not establish perceived comfort
through a reported participant study. This leaves room for later HRI
validation, but not a claim that comfort was wholly absent from prior work.
\end{itemize}

These observations are not a re-evaluation of the paper's code or raw data.
They identify which claims its printed evidence can support and which
quantities would need harmonization in a comparative experiment.

Omni-Perception is another close 2025 comparator: its abstract describes raw
LiDAR, spatiotemporal risk processing, omnidirectional avoidance, and a
simulation toolkit compatible with multiple engines \citep{omni}. It prevents
novelty from being assigned merely to direct point-cloud perception or
multi-simulator raycasting. Its detailed comparative results were not
independently inspected in this review.

\subsection{FLUX: a 2026 cross-embodiment navigation competitor}
FLUX uses rectified-flow trajectory generation followed by dynamic-scene
reinforcement-learning refinement, and introduces DynBench \citep{flux}.
The inspected version includes static and dynamic navigation tasks and
reports transfer to wheeled, quadrupedal, and G1 humanoid platforms. Its
deployment uses an RGB-D camera and inference on a remote RTX~4060 machine.
This matters when comparing onboard compute constraints and communication
latency with Geoff-Sim.

Three details are especially relevant to critical positioning:
\begin{enumerate}
\item Its Table~II reports SocialNav success of 64.0 versus 59.3 for NavDP.
This is meaningful within that protocol, but goals, episode termination,
pedestrian speeds, and social-distance definitions differ from Geoff-Sim.
It is not a directly comparable headline improvement.
\item Section~V defines ``Collision Count'' as intimate-space intrusions
below 0.45\,m, a proxy for collision risk. Later prose describes values
with percent signs while the metric definition describes counts.
Accordingly this review does not reinterpret those values as physical
contact probabilities. Physical contacts and proxemic intrusions should
be distinct outputs in the proposed study.
\item Table~IV reports 216.5\,ms inference for the selected six-step
configuration. This is a different timing boundary and compute setup from
the project's local-planner p95 measurements. It motivates measuring
end-to-end observation-to-actuation age; it does not prove one system
is faster under matched conditions.
\end{enumerate}

FLUX's limitation discussion acknowledges missing explicit social reasoning
and predominantly simulated indoor pedestrian evaluations. Its existence
still decisively shifts the research gap: cross-embodiment dynamic
navigation is no longer an untouched area. A narrowly scoped G1 response
and uncertainty study can be complementary to such general policies.

\section{Safety, middleware, and whole-system timing}
\label{sec:safety}
\subsection{Filtering commands versus guaranteeing safety}
Control barrier functions provide a formal route to invariant-set conditions
and quadratic-program safety filtering under stated model assumptions
\citep{cbf}. The safe-learning survey distinguishes empirical robustness,
learning-based improvement, and formal certification \citep{safecontrol}.
These distinctions should govern terminology in this project. A collision
supervisor that checks predicted rollouts is a safety-oriented mechanism;
without validated disturbance bounds and recursive feasibility, it is not
automatically a safety certificate.

A navigation command $u$ is transformed before its physical effect appears.
Let $\mathcal{N}$ denote normalization, deadbands, and command shaping, and
let $\hat f_\theta$ approximate policy-level motion response. A candidate
rollout can be expressed as
\begin{equation}
\hat x_{k+1}=\hat f_\theta(\hat x_k,\mathcal{N}(u_k),d_k),
\end{equation}
where $d_k$ includes the measured or modeled delay. This notation summarizes
a proposed experimental model, not a claimed full-order description of G1.
Its accuracy must be measured during braking, turning, and combined-axis
motion, not only steady straight walking.

For a one-dimensional stop with constant speed during delay and constant
deceleration afterward, an illustrative distance is
\begin{equation}
d_{\rm stop}=v d+\frac{v^2}{2a_{\rm brake}},\qquad a_{\rm brake}>0.
\label{eq:stop}
\end{equation}
The equation explains why zero commanded velocity is not zero stopping
distance. It is not a bound until delay and achievable braking have been
validated, and it omits a moving pedestrian's continuing approach.
For example, if a person approaches at 1\,m/s, an observation age of
0.305\,s corresponds to 0.305\,m of unobserved motion before future-horizon
prediction even begins. This is an illustrative calculation using a chosen
human speed, not a reconstructed collision from the repository.

\subsection{Planner--supervisor interaction}
A planner may choose a trajectory that the supervisor repeatedly vetoes.
Even when every individual veto seems prudent, repeated disagreement can
cause oscillation or starvation. A suitable study measures override rate,
rejected-command reasons, clearance under the actual applied command,
yield duration, and progress after resumption. It should also distinguish
stopping in front of an approaching person from moving to a genuinely
lower-risk configuration: stationary is not synonymous with safe in a
dynamic environment.

For a humanoid, emergency velocity changes can create a second hazard:
instability or contact by a swinging body segment. A planar footprint can
be useful if conservatively validated, but it does not directly measure
full-body swept volume, fall reach, or contact force. Such distinctions are
particularly relevant to the project's use of independent pedestrian actors.

\subsection{ROS~2, Nav2, and SLAM are enabling infrastructure}
ROS~2 provides the communication and lifecycle architecture \citep{ros2}.
Navigation2's behavior-tree orchestration is described in The Marathon~2
\citep{nav2}, and SLAM Toolbox supports the mapping component \citep{slam}.
MuJoCo provides the physics foundation \citep{mujoco}. Combining these tools
is valuable research engineering, but their presence is not a new navigation
algorithm. Nor does middleware design establish hard real-time performance
for an entire application.

The current event-driven route monitor, forecast costmap layer, and retained
path execution offer a specific timing experiment: when does event-triggered
replanning improve collision-free progress relative to fixed-rate replanning
under the same compute budget? A 5\,Hz request limit is not a 200\,ms route
delivery guarantee. Queueing, costmap availability, retries, and action
transport belong in the latency measurement. Global predictive costs and
local collision checking should also be ablated independently to avoid
crediting the same forecast twice.

\section{Benchmarking and social validity}
\label{sec:benchmarks}
\subsection{Existing benchmarks already cover important parts of the problem}
SocNavBench provides grounded scenarios, consistent metrics, and a common
testing framework \citep{socnavbench}. SocialGym~2.0 emphasizes configurable
multi-robot navigation in shared human spaces \citep{socialgym}. Habitat~3.0
adds human avatars, human-in-the-loop interaction, and collaborative tasks
\citep{habitat3}. MetaUrban extends compositional urban environments and
heterogeneous mobility evaluation \citep{metaurban}. These address realism,
comparability, interaction, or scale that a small custom plaza cannot assume.

OmniNavBench is particularly important to a 2026 gap analysis
\citep{omninav}. Its inspected version includes 170 environments and 1,779
human-teleoperated trajectories, multiple embodiments including Unitree H1,
continuous physical stepping, and composed navigation tasks. The low
complete-task success of several evaluated models highlights a distinction
between individual skills and reliable task sequences. Its social metric
uses a personal-space threshold, which should not be equated with a
full-body contact sensor. Geoff-Sim's proposed advantage should therefore
be fine-grained control/measurement provenance and calibrated execution
analysis, rather than the first physically embodied multi-task benchmark.

The very recent DPed-VLN abstract describes language navigation with dynamic
pedestrians, paired instructions, and social-safety metrics \citep{dped}.
It is adjacent rather than a direct baseline for a LiDAR/Nav2 waypoint stack.
Its inclusion demonstrates why ``navigation in dynamic pedestrian scenes''
is too broad a novelty claim. This review does not infer its low-level
humanoid dynamics or deployment properties from the presence of humanoid
pedestrian avatars.

\subsection{Safety is necessary but not sufficient for social navigation}
The evaluation guidelines articulate safety, comfort, legibility,
politeness, social competency, agent understanding, proactivity, and context
responsiveness \citep{guidelines}. The 2026 proxemics taxonomy further
emphasizes human, robot, environment, and contextual influences
\citep{proxemics2026}. A fixed circular personal-space penalty is consequently
a controllable proxy, not a universal measure of comfort.

A simulator can validly measure minimum body clearance, contact events,
time in specified proxemic regions, approach speed, jerk, and route progress.
To infer perceived safety or comfort, the study needs appropriately designed
human judgments, such as counterbalanced video ratings or later controlled
in-person comparisons. These should separate morphology, speed, yielding
style, and context. A humanoid's natural-looking gait may improve legibility
or may increase expectations that the controller does not meet; the direction
of the effect is an empirical question.

\subsection{A useful validity hierarchy}
\begin{enumerate}
\item \textbf{Mock regression:} verifies scenario, scoring, aggregation, and
failure-handling mechanics; it provides no physical avoidance evidence.
\item \textbf{Oracle physics:} tests planning and execution with privileged
perception; it cannot establish sensor-only performance.
\item \textbf{Onboard-sensor simulation:} tests the integrated observation and
execution pipeline under the simulator's sensing/physics assumptions.
\item \textbf{Cross-simulator or perturbed-model evaluation:} probes sensitivity
to selected modeling assumptions; it is not itself hardware transfer.
\item \textbf{Hardware and human evaluation:} addresses physical transfer and
human-centered outcomes under an explicit deployment protocol.
\end{enumerate}
The current project spans the first three levels with incomplete reliability
at level three. Reporting the hierarchy avoids overstating either the mock
tests or the passing oracle demonstrations.

\section{Comparative synthesis and ranked research gaps}
\label{sec:gaps}
\begin{table}[htbp]
\small
\caption{Closest families and the remaining project-specific question.
This is a qualitative synthesis, not a benchmark ranking.}
\begin{tabularx}{\textwidth}{P{0.23\textwidth}YY}
\toprule
Prior work & Established capability & Question still testable here \\
\midrule
DWA, VO, TEB, biped crowd navigation & Dynamics/feasibility matter in avoidance & Does a calibrated policy-level response envelope outperform generic bounds? \\
CADRL, SARL, RGL, DS-RNN & Learned spatial/temporal crowd reasoning & What happens under matched sensing age, misses, and G1 execution? \\
ASAP, RMA, perceptive locomotion & Adaptation and dynamics-aware learned control & Can frozen-policy response calibration improve navigation without retraining locomotion? \\
Safe humanoid policy, Omni-Perception & Sensor-conditioned legged/humanoid avoidance & How do modular filtering and end-to-end policies trade contact risk against liveness? \\
FLUX, OmniNavBench & Cross-embodiment navigation and richer evaluation & How much apparent performance changes with information provenance and measured execution delay? \\
Social-navigation guidelines and proxemics taxonomy & Multidimensional human-centered evaluation & Which safety/progress gains correspond to judged comfort and legibility? \\
\bottomrule
\end{tabularx}
\end{table}

\subsection{Gap 1: calibrated executable avoidance under stale observations}
\textbf{Priority: highest; fit: direct; novelty confidence: moderate and
conditional.} Established feasible-velocity methods, biped planning, and
recent LiDAR policies already address parts of the problem
\citep{vo,dwa,biped2016,safehumanoid}. The narrower unresolved question for
this platform is the contribution of a \emph{measured policy-level response
distribution} combined with observation-age-aware person occupancy.

\textbf{Hypothesis H1.} At matched clearance conservatism and compute budget,
a calibrated response envelope lowers contact incidence or increases
collision-free mission completion compared with instantaneous and
hand-tuned response models.

\textbf{Contribution required.} Define what is calibrated: delay, gain,
acceleration/braking asymmetry, yaw response, combined-axis coupling, and
residual prediction error. Condition or stratify by relevant initial motion.
Use held-out commands and scenes, and propagate residual uncertainty into
the navigation decision. The estimator and its uncertainty treatment must
be specified sufficiently to reproduce.

\textbf{Falsification.} If an equally conservative fixed-margin baseline
matches the result, or improvements disappear on held-out response tests,
the extra model complexity is not justified. A finding that timing dominates
model complexity would still be useful if established through controlled
interventions.

\subsection{Gap 2: causal explanation of the oracle-to-onboard gap}
\textbf{Priority: high; fit: direct; novelty confidence: moderate.}
Benchmarks and perception datasets already exist
\citep{socnavbench,jrdb,omninav}. What can be distinctive is a source-audited
experimental decomposition on the \emph{same walking system}.

\textbf{Hypothesis H2.} During close encounters, track age and partial-body
velocity bias explain more avoidable failures than map or localization
error alone, after controlling for scenario and locomotion policy.
This is a hypothesis suggested by the repository's diagnostics, not an
established causal result.

\textbf{Contribution required.} Independently replace observed person tracks,
estimated robot pose, and online map with their oracle counterparts. Keep
goal definitions, actor motion, policy, and evaluation truth fixed wherever
possible. Then inject controlled delay, track dropout, identity changes,
and pose drift. Log which information was actually consumed at the source
selection branch, rather than relying only on launch arguments.

\textbf{Falsification.} If failures persist under oracle people but disappear
under a changed low-level policy, the principal cause is likely execution
rather than perception. If the accepted goals or actor exposure change
between sensing modes, the comparison is confounded and needs redesign.

\subsection{Gap 3: coordinated safety and liveness across planning layers}
\textbf{Priority: high; fit: direct; novelty confidence: moderate-to-low
until a specific decision rule is defined.} Freezing, cooperation, and
hierarchical navigation are already studied
\citep{freezing,rgl,parkour,wheeledleg}. The opening is not the existence of
a global/local hierarchy, but a verified policy for resolving disagreement
between route intent, social planning, and safety supervision.

\textbf{Hypothesis H3.} Response-consistent handoff and event-driven route
updates reduce prolonged veto loops and timeout probability without
increasing contact or fall rates compared with fixed-rate replanning and
independent supervision.

\textbf{Contribution required.} Define the state machine or optimization
criterion governing yield, resume, detour, and route replacement. Measure
delivered-route age, discarded or stale routes, supervisor disagreement,
recovery time, and progress after a yield. Match total CPU budget and
prediction assumptions across alternatives.

\textbf{Falsification.} Higher route-update frequency without improved
collision-free completion is an engineering change, not confirmation of
the hypothesis. Likewise, reduced timeouts bought by more contacts do not
establish a better trade-off unless the full frontier is reported.

\subsection{Gap 4: full-body risk beyond planar clearance}
\textbf{Priority: medium; fit: requires extensions; novelty confidence:
conditional.} Whole-body control and three-dimensional avoidance already
exist \citep{expressive,asap,safehumanoid,omni}. A useful narrower question
is when a planar swept footprint is inadequate for humanoid social navigation.

\textbf{Hypothesis H4.} A validated body-swept-volume approximation detects
encounter risks missed by base clearance alone, especially during turning,
lateral stepping, and recovery.

This requires link-level contact/clearance reference measurements and
scenario designs that exercise torso/arm motion and stability. Start by
quantifying the error of the current footprint abstraction. A more expensive
full-body method is warranted only if that error affects decisions. Claims
about contact injury or human safety would need a different validation
scope than collision counting.

\subsection{Gap 5: contextual comfort and nonreciprocal crowds}
\textbf{Priority: longer term; fit: requires human evaluation and richer
actors.} Social principles and contextual proxemics are well developed
\citep{guidelines,proxemics2026}. The opportunity is to test whether
response-aware motion creates behavior people actually find more
predictable and comfortable, while avoiding an assumed universal radius.

Compare nonreactive, reciprocal, and data-driven pedestrian models first.
Then evaluate selected robot behaviors with human judgments, controlling
for speed, clearance, and morphology. Without this step, use the term
\emph{human-aware collision avoidance} or \emph{proxemic-cost navigation}
for the demonstrated capability rather than claiming broadly validated
social intelligence.

\section{A concrete experimental programme}
\label{sec:experiments}
\subsection{Stage A: characterize the frozen locomotion interface}
Collect command step, ramp, stop, turn, and combined-axis responses across
initial speeds, directions, and repeated runs. Record requested command,
normalized command, supervisor output, policy-consumed command, estimated
state, evaluation state, and timestamps. Measure latency, gain, overshoot,
stopping distance, lateral drift, and fall outcomes. Separate training,
calibration, and held-out response sets by complete trajectories, not
adjacent samples from one rollout.

Compare (i) instantaneous planar integration; (ii) current hand-set envelope;
(iii) fitted deterministic response; and (iv) fitted response with residual
uncertainty. Evaluate horizon-dependent error and interval coverage. Software
slew limits are not automatically physical acceleration bounds. Calibrate
the active profile and policy together; switching either changes the model.

\subsection{Stage B: controlled oracle planning study}
Freeze maps, source information, policy, robot geometry, actor trajectories,
and goals. The repository's development methods---Nav2-only, reactive
geometric, deterministic temporal, and full planner/supervisor---provide a
starting matrix, but are not substitutes for published comparator
implementations. Add a well-specified classical baseline, preferably DWA,
VO/ORCA with explicit reciprocity settings, or TEB, and describe adaptation
to the G1's command interface.

All methods must share unavoidable hardware/policy limits. If the same
supervisor is applied to every method, state that the comparison concerns
the upstream planner under a shared shield. If supervision itself is the
independent variable, separate the experimental arms so that a shared
filter does not erase the difference being measured. Match tuning budgets;
select parameters on development scenarios only.

\subsection{Stage C: factorial onboard robustness study}
Use a small, interpretable factorial design before scaling to many methods.
Suggested factors are response model, observation-age compensation,
perception source, and forecast uncertainty treatment. Add localization/map
substitutions as diagnostic blocks if the full factorial is too costly.
Report interactions: for example, a response model may help only when
track age is corrected.

Scenario families should include head-on passing, orthogonal crossing,
overtaking, occlusion emergence, corridor blockage, narrow passage, moving
crowds, and turn/stop transients. Vary crowd density, human-to-robot speed
ratio, sensor dropout, map ambiguity, and actor reciprocity. Include
people-free trials to isolate locomotion and route-following failures.
Hold out at least one scene family to probe overfitting to the plaza.

Identical seeds are not sufficient for pairing live runs: wall-time
scheduling and observation-dependent goal admission can change encounters.
For causal replay, freeze actor trajectories and goals where feasible,
record accepted events, and quantify residual nondeterminism. Treat
nonreactive replay and interactive crowds as separate experimental regimes.

\subsection{Stage D: metrics and statistical interpretation}
Predeclare primary outcomes: complete mission without contact/fall, contact
incidence, and algorithmic timeout. Secondary outcomes include minimum
body clearance, contact count/severity proxy, time in specified proxemic
regions, approach speed, stopping error, path/time efficiency, yield/recovery
duration, supervisor interventions, and processing/observation/actuation
latencies. State whether a contact metric is based on geometry, physics
contacts, or a distance threshold.

Keep four denominators visible: planned trials, executed trials, valid
trials, and successful trials. A collision is an algorithmic outcome; a
missing required scoring stream is an experiment-validity problem. However,
method-induced overload can itself cause missing data. Predeclaring failure
classification and reporting invalidity by method prevents invalid-trial
exclusions from hiding poor computational robustness.

Use confidence intervals for rates and paired scenario/seed differences
where replay supports pairing. Bootstrap at the independent scenario/seed
unit, with repetitions nested within it, not at individual telemetry samples.
The repository's Wilson intervals and paired aggregation are useful
starting points, but the statistical design still needs to specify the
independent unit and missingness assumptions. For zero observed collisions
in $n$ independent trials, the one-sided exact 95\% upper bound is
\begin{equation}
p_{\rm upper}=1-0.05^{1/n}.
\end{equation}
Even zero events in 100 independent trials leaves an upper bound near 3\%.
Repeated correlated scenarios do not justify the same interpretation.
Choose the trial count from desired precision or a prespecified detectable
effect, not from the first sequence of passing demonstrations.

Do not rank methods only by travel time on successful trials: this conditions
on success and can reward a fast but unreliable method. Report completion
and failure distributions alongside conditional efficiency. Avoid a single
weighted score obscuring collisions, falls, invalidity, and timeouts.

\subsection{Minimum publishable evidence package}
\begin{enumerate}
\item A frozen protocol, policy/configuration hashes, source manifests,
scenario definitions, and explicit information boundaries.
\item Held-out response-model error and calibration results.
\item Matched baselines plus ablations isolating response, delay,
uncertainty, and supervision.
\item Multi-scenario, multi-seed mission results with retained failures,
validity accounting, and uncertainty estimates.
\item Failure-case analysis linking sensor observations, decisions, applied
commands, and physical outcomes.
\item Reproducible artifacts and a bounded contribution statement supported
by those experiments.
\end{enumerate}

\section{Recommended positioning and conclusion}
\textbf{Suggested working title:}
\emph{Observation-Age-Aware, Response-Calibrated Social Navigation for a
Walking Humanoid: A Source-Audited Evaluation in MuJoCo and ROS~2.}

This title describes a proposed research direction, not completed results.
The paper's eventual central claim should take a falsifiable form:
\begin{quote}
Under a declared sensing, locomotion, and crowd model, the proposed
calibration and decision procedure improves collision-free mission
completion relative to matched baselines, and controlled ablations identify
when the improvement depends on observation age and executable response.
\end{quote}
Whether that statement becomes true is an experimental question. At present,
the repository supports the platform and diagnosis portions of the work,
while reliable onboard avoidance remains open.

The literature gives a clear boundary to novelty. Dynamics-aware avoidance,
biped crowd planning, learned crowd interaction, sensor-driven G1 avoidance,
cross-embodiment navigation, and physically grounded navigation benchmarks
already exist. The most promising contribution is a carefully controlled
account of their integration failure modes, together with a specific,
validated mechanism that improves the safety--liveness trade-off. A rigorous
negative result showing that a proposed response model does not help under
realistic sensing would be more informative than an unqualified success
claim based on selected oracle runs.

\appendix
\section{Annotated evidence map of the 50 retained works}
\label{app:evidence}
The entries below give each work's role in this review. They are not claims
that every work received an equally detailed full-text review. Identity
records and content-inspection levels are listed separately in the audit.

\small
\begin{longtable}{P{0.28\textwidth}P{0.65\textwidth}}
\toprule
Work & Contribution to the synthesis and transfer limitation \\
\midrule
\endfirsthead
\toprule Work & Contribution to the synthesis and transfer limitation \\
\midrule
\endhead
\bottomrule
\endfoot
Kruse et al. \citep{survey2013} & Foundational human-aware navigation framing; supplies historical context rather than a contemporary G1 baseline. \\
Singamaneni et al. \citep{survey2024} & Faceted taxonomy of socially aware navigation; guards against reducing social behavior to clearance alone. \\
Learning-based review \citep{survey2023} & Recent comparative overview and benchmark perspective; motivates learned baselines without transferring its rankings to different sensing contracts. \\
Evaluation guidelines \citep{guidelines} & Defines multiple social-navigation principles and evaluation dimensions; supports separate safety and human-centered outcomes. \\
Nahum et al. \citep{proxemics2026} & Context-dependent proxemics taxonomy; argues against universal personal-space geometry as a complete comfort metric. \\
MuJoCo \citep{mujoco} & Physics-engine foundation; simulator choice alone is neither calibrated dynamics nor algorithmic novelty. \\
ROS 2 \citep{ros2} & Middleware architecture and deployment context; application-level latency remains a measured property. \\
The Marathon 2 \citep{nav2} & Nav2 architecture and navigation orchestration; not a G1 policy-response model. \\
SLAM Toolbox \citep{slam} & Mapping infrastructure; online mapping success must be separated from crowd avoidance. \\
DWA \citep{dwa} & Dynamics-constrained velocity search baseline; makes generic acceleration-aware novelty claims inappropriate. \\
Velocity obstacles \citep{vo} & Relative-motion collision reasoning and admissible velocities; requires appropriate motion estimates and robot constraints. \\
ORCA \citep{orca} & Reciprocal collision-avoidance reference; actor reciprocity must match or be explicitly varied. \\
Sparse trajectory optimization \citep{teb} & Time/trajectory optimization comparator; the deployed controller model remains consequential. \\
Social force \citep{socialforce} & Interpretable pedestrian interaction model; not universal empirical human behavior. \\
Trautman et al. \citep{freezing} & Dense-crowd cooperation perspective; motivates explicit freezing/progress metrics. \\
CADRL \citep{cadrl} & Learned decentralized collision avoidance; representation and training crowd assumptions need matching. \\
SARL \citep{sarl} & Attention-based crowd interaction reasoning; useful learned-state baseline, separate from person detection. \\
RGL \citep{rgl} & Graph relations and predictive planning; compare under the same sensing and execution model. \\
DS-RNN \citep{dsrnn} & Temporal/spatial learning and partial observability; real wheeled transfer does not establish biped response equivalence. \\
SocNavBench \citep{socnavbench} & Grounded, common evaluation scenarios and metrics; already occupies broad benchmark-novelty territory. \\
SocialGym 2.0 \citep{socialgym} & Configurable multi-agent navigation environments; robot interaction assumptions affect transfer. \\
JRDB \citep{jrdb} & Egocentric human perception dataset; relevant to misses/association, not a complete G1 closed-loop benchmark. \\
SCAND \citep{scand} & Human navigation demonstrations across robot morphologies; behavioral reference with different sensors and actuation. \\
DROW \citep{drow} & Learned sparse-range detection precedent; wheelchair-oriented original task differs from torso-mounted G1 person tracking. \\
DR-SPAAM \citep{drspaam} & Temporal/spatial person detection from range scans; projection and walking-motion domain shift require evaluation. \\
FAST-LIO2 \citep{fastlio} & Strong direct LiDAR--inertial estimation comparator; localization improvements alone do not establish navigation reliability. \\
ORB-SLAM3 \citep{orbslam} & Visual/visual--inertial state-estimation alternative; requires a suitable camera pipeline. \\
Social LSTM \citep{sociallstm} & Interaction-aware learned trajectory prediction; forecasting and decision quality are different outcomes. \\
Social GAN \citep{socialgan} & Diverse plausible futures; sample diversity does not by itself calibrate risk. \\
Trajectron++ \citep{trajectron} & Dynamical, heterogeneous, optionally ego-conditioned prediction; an important alternative to independent constant velocity. \\
Rudenko et al. \citep{prediction} & Forecasting taxonomy, context, datasets, and metrics; motivates matching the prediction task to planning. \\
CBF-QP \citep{cbf} & Formal safety filtering under assumptions; distinguishes invariant conditions from empirical rollout checks. \\
Safe learning survey \citep{safecontrol} & Separates safe control, robustness, and certification; prevents overclaiming trained or heuristic safety. \\
Biped crowd navigation \citep{biped2016} & Direct historical prior art on biped safety in crowds; detailed experiments were not inspected because of access limitations. \\
RMA \citep{rma} & Rapid locomotion adaptation; adaptation competence is not a pedestrian-planning policy. \\
Massively parallel walking \citep{rudin} & Efficient locomotion training; does not establish social mission completion. \\
Perceptive locomotion \citep{rough} & Robust exteroceptive/proprioceptive integration; primarily terrain competence rather than social acceptability. \\
ANYmal Parkour \citep{parkour} & Capability-aware navigation hierarchy; demonstrates that embodiment-aware planning already exists. \\
Wheeled-legged autonomy \citep{wheeledleg} & Integrated mobility-aware navigation and long-range missions; different morphology but strong systems-level prior art. \\
ExBody \citep{expressive} & Expressive humanoid control with locomotion; perceived comfort needs an additional evaluation. \\
DeepMimic \citep{deepmimic} & Physics-based motion imitation foundation; animation/skill objectives differ from crowd safety. \\
Dynamics randomization \citep{domain} & Foundational sim-to-real robustness method; not a calibrated stopping certificate. \\
ASAP \citep{asap} & Residual action alignment and G1 transfer; full-text methods/tables examined, task is motion imitation rather than social navigation. \\
Safe humanoid locomotion \citep{safehumanoid} & Closest LiDAR/G1 safety/comfort competitor; inspected costs and tables support a qualified, scenario-dependent comparison. \\
FLUX \citep{flux} & Recent generative navigation and G1 transfer; inspected metric definitions and deployment timing constrain cross-paper comparisons. \\
OmniNavBench \citep{omninav} & Physically embodied, multi-task benchmark; inspected protocol narrows the open benchmark contribution. \\
MetaUrban \citep{metaurban} & Compositional urban micromobility simulation; abstract-screened broader environment/embodiment coverage. \\
Habitat 3.0 \citep{habitat3} & Human/avatar collaboration and human-in-loop evaluation; abstract-screened scope differs from low-level G1 response. \\
DPed-VLN \citep{dped} & September 2026 dynamic-pedestrian language-navigation benchmark; abstract-only emerging evidence, not a direct Nav2 comparator. \\
Omni-Perception \citep{omni} & Raw-LiDAR omnidirectional legged avoidance; abstract-screened close competitor requiring deeper comparison before publication. \\
\end{longtable}
\normalsize

\section{Claim audit and immediate priorities}
\begin{longtable}{P{0.43\textwidth}P{0.50\textwidth}}
\toprule
Potential claim & Assessment and required evidence \\
\midrule
\endfirsthead
\toprule Potential claim & Assessment and required evidence \\
\midrule
\endhead
\bottomrule
\endfoot
``First humanoid social-navigation system'' & Contradicted by direct prior art; do not use. \\
``First dynamics-aware avoidance'' & Contradicted by DWA, VO, biped planning, and subsequent embodied methods. \\
``Novel response-calibrated uncertainty-aware method'' & Possible only after specifying the new mechanism and comparing the nearest methods under matched conditions. \\
``Sensor-only navigation implemented'' & Supported by the documented information boundary and tests, within its stated graph/configuration scope. \\
``Reliable onboard avoidance achieved'' & Not supported by the current documented collisions and timeouts. \\
``Hard real-time planning'' & Not supported by a soft search budget or observed p95 runtime. \\
``Formal safety guarantee'' & Requires valid model/disturbance assumptions and an appropriate proof; heuristic rollout filtering is insufficient. \\
``Socially comfortable motion'' & Geometric proxies may be reported as proxies; perceived comfort requires human-centered evidence. \\
``Reproducible benchmark contribution'' & Plausible with validated artifacts, controlled replay, external baselines, and new scientific findings; seeded development demos alone are insufficient. \\
``Simulation results transfer to hardware'' & Requires hardware evidence; neither MuJoCo realism nor another paper's G1 deployment establishes this project's transfer. \\
\end{longtable}

\clearpage
\bibliographystyle{plainnat}
\bibliography{references}
\end{document}
